\section{Conclusion}\label{sec:conclusion}

This paper introduced the ZISAR (Zero-Inflated Spatial Autoregressive) model, a unified specification for count data that simultaneously addresses spatial spillovers, zero-inflation, overdispersion, and unobserved unit-level heterogeneity.
The key innovation is the placement of a simultaneous autoregressive structure on the latent linear predictor of the count process---$\boldsymbol{\eta} = (\mathbf{I} - \rho\bW)^{-1}(\bX\bbeta + \bu)$---rather than on the observed count.
This SAR-on-$\boldsymbol{\eta}$ formulation preserves the direct/indirect/total effects decomposition that is central to the spatial econometrics programme \citep{lesage_pace2009,arbia2006} while avoiding the endogeneity of observed spatial lags that complicates SAR-on-$\by$ specifications \citep{lambert2010}.
Simultaneously, it differs from CAR-based spatial random-effects models \citep{lee2013,neelon2013} by providing a structural, single-parameter characterisation of spatial dependence with an explicit spillover interpretation.

Estimation by Laplace-approximated marginal likelihood, implemented via \texttt{TMB} \citep{kristensen2016}, yields a computationally efficient and scalable algorithm.
The sparse-matrix operations and automatic differentiation exploited by \texttt{TMB} echo the computational concerns raised by \citet{arbia_nardelli2024} and \citet{arbia_nardelli_salvini} in the context of pairwise pseudo-likelihood for large spatial datasets.
We established consistency and asymptotic normality of the estimator and derived the effects decomposition for the nonlinear zero-inflated mixture.

The Monte~Carlo study confirmed that ignoring either spatial dependence or zero-inflation leads to substantial bias: the standard ZINB inflates covariate effects when $\rho > 0$, while the SAR-NB severely underestimates the spatial lag when structural zeros are present, with a bias that does not vanish as $n$ grows.
ZISAR recovers the data-generating parameters with low bias and an RMSE that decreases with $n$.
A second experiment showed that residual-based Moran tests can miss spatial dependence that the likelihood detects.

The application took the estimator to the $n = 32{,}322$ small areas of England, where a fit takes five to six minutes on a laptop.
On domestic solar PV installations during the boom of the Feed-in Tariff, before the tariff cut of 2016, it needs both parts of the restricted model.
Adoption is strongly dependent in space ($\hat\rho = 0.63$, likelihood ratio 8{,}463 against the ZINB with random effects), and zero-inflation remains needed once the spatial field is in (likelihood ratio 552 against the SAR-NB-RE), so that ZISAR has the lowest AIC and BIC of six specifications.
The structural zeros follow the housing stock: nearly half of the observed zeros are structural, three quarters of them are in Greater London, and their posterior probabilities are higher where areas remained without installations after the cut, which the model does not use.
The aspatial models read the common clustering of tenure, property value and adoption as effects of the covariates, which shrink by an order of magnitude in ZISAR, and they have no counterpart to the indirect effects, which are about 60\% of every total effect.
The window is chosen so that zeros are informative: during the boom so many installations are expected in most areas that a zero is improbable under the count process, whereas after the cut, with one or two installations per area, a spatially dependent count process reproduces clusters of zeros as sampling zeros.
The same high intensities make the Laplace objective non-smooth, and the optimiser stops with a large gradient at an optimum that a simulation from the fitted model shows to be genuine; a parametric bootstrap provides the standard errors.

The general model, with a spatial field in the inflation predictor, is not needed in this application, since the covariates explain where the structural zeros are.
Where it is, the simulations of Appendix~\ref{app:general} show that it should be estimated with a reduced-rank field on a lattice of knots, which the Laplace approximation handles, rather than with a field of one element per unit, which produces likelihood-ratio statistics in the thousands on data that contain no field at all.

\paragraph{Limitations.}
Four limitations warrant acknowledgement.
First, the Laplace approximation may lose accuracy when the conditional posterior of $\boldsymbol{\eta}$ is highly non-Gaussian (e.g.\ for very sparse data with many zero counts); higher-order Laplace corrections or adaptive quadrature could improve precision in such settings.
With the unit-level field of the general model it produced spurious optima, with gradients of order $10^8$ at the stopping point and log-likelihoods thousands of points above their true value (Appendix~\ref{app:general}).
A large gradient at the stopping point flags them but is not specific: with structural zeros among units of high intensity, as in the solar application, the conditional posterior of $\eta_i$ of a zero unit has two modes and the optimiser stops with a large gradient at a genuine optimum, so that a simulation from the fitted model is needed to tell the two cases apart, and a parametric bootstrap replaces the Hessian standard errors.
Second, the unit-level random effects $\bu$ are assumed independent Gaussian; correlated or non-Gaussian specifications would add flexibility at the cost of a more complex \texttt{TMB} template.
Third, the eigenvalue pre-computation for the log-determinant $\log|\det(\mathbf{I}-\rho\bW)|$ is $O(n^3)$; the sparse-Cholesky alternative of \eqref{eq:logdet_sparse} removes this bottleneck at the national scale of the application, but its cost grows with the fill-in of the Cholesky factor, so that for dense weights matrices or $n$ in the millions stochastic log-determinant approximations \citep[e.g.][]{barry_pace1999} would be needed.
Fourth, the general specification with a unit-level field in the inflation predictor is poorly served by the Laplace approximation: each element of the field is informed by a single outcome that is zero or not, the approximation error accumulates over the units, and on data generated without a field the approximate likelihood-ratio statistic runs into the thousands.
The reduced-rank field on a lattice of knots removes the failure and is estimated without difficulty on simulated data at the national scale, at the price of a resolution set by the knot spacing: a field whose variance lies below the spacing is detected but not reconstructed (Appendix~\ref{app:general}), and as the spacing falls towards the size of the units the unit-level case and its failure return.
Simulation-based inference on the same \texttt{TMB} template, such as Hamiltonian Monte Carlo, remains the route to a unit-level field.

\paragraph{Extensions.}
Several directions for future research emerge naturally.
The most immediate is the incorporation of \emph{temporal dependence}, yielding a spatio-temporal ZISAR of the form
\[
  \eta_{it} = \rho(\bW\boldsymbol{\eta}_t)_i + \phi\,\eta_{i,t-1} + \mathbf{x}_{it}'\bbeta + u_i + \nu_t,
\]
where $\phi$ captures persistence and $\nu_t$ is a time-specific effect.
The panel structure introduces additional identification and computational challenges (dynamic spatial panel asymptotics, high-dimensional latent field) that merit a dedicated treatment.

A second extension concerns alternative spatial structures: replacing the single $\bW$ with multiple weight matrices (e.g.\ geographic and economic proximity), or letting $\bW$ be endogenous.

Third, the ZISAR template naturally accommodates other count distributions (generalised Poisson, Conway--Maxwell--Poisson) and alternative inflation links (probit, complementary log-log) by modifying the observation-model block of the \texttt{TMB} template.

Finally, automated tools for choosing among the restricted and general specifications, and among weights matrices, would make the \texttt{zisar} package \citep{zisar_pkg} easier to use in applied work.
